\documentclass[aspectratio=169]{beamer}
\usetheme{Madrid}
\usecolortheme{default}

\usepackage{amsmath}
\usepackage{booktabs}
\usepackage{graphicx}
\usepackage{siunitx}
\usepackage{xcolor}

\title[PDC Reconstruction and Uncertainty]{Detailed PDC Target-Momentum Reconstruction and Uncertainty Analysis}
\subtitle{Methods, Error Models, and Confidence/Credible Intervals}
\author{SMSimulator5.5 Analysis}
\date{\today}

\newcommand{\LargeRunRoot}{../../../results/pdc_rk_error_after_pdcana_20260411}
\newcommand{\SmokeRunRoot}{../../../results/pdc_rk_error_after_pdcana_smoke3}
\newcommand{\PlotOrFallback}[3]{%
  \IfFileExists{#1}{\includegraphics[width=#3\textwidth]{#1}}{%
    \IfFileExists{#2}{\includegraphics[width=#3\textwidth]{#2}}{%
      \fbox{\parbox[c][0.45\textheight][c]{0.80\textwidth}{\centering Plot not available yet}}}}}

\begin{document}

\begin{frame}
\titlepage
\end{frame}

\begin{frame}{Coordinate frames: lab vs.\ target (beam-as-Z)}
\small
Two right-handed frames:
\begin{itemize}
  \item \textbf{Lab frame} (Geant4 world): $+z$ along the nominal SAMURAI beamline.
    PDC hits, field map, RK integration, and the internal fit live here.
  \item \textbf{Target frame} (``beam-as-Z''): $+z$ along the incident beam;
    rotated about $y$ by $-\alpha_\mathrm{tgt}=-3^\circ$ relative to lab.
    \texttt{truth\_proton\_p4} stays in this frame
    (see \texttt{DeutPrimaryGeneratorAction.cc}: the forward $R_y(-\alpha)$ is
    applied to a local copy only; \texttt{fMomentum} is never rewritten).
\end{itemize}

\textbf{All truth vs.\ reco plots and all interval coverage numbers in this
report are in the target frame.} Reco tools apply $R_y(+\alpha_\mathrm{tgt})$
at the write stage; the $3\times3$ momentum covariance is rotated as
$\Sigma'=R\Sigma R^\mathrm{T}$; $px/py/pz$ Fisher/Laplace intervals are
Gaussian-rebuilt from the rotated diagonals; $|\vec p|$ is rotation-invariant.
MCMC/profile asymmetric quantile information degrades to Gaussian fallback.
Prior to 2026-04-19 this rotation was missing; the affected
$\Delta p_x \approx -p_z\sin\alpha \approx -33$~MeV/$c$ bias contaminated
all pre-fix plots ($p_z \approx 627$~MeV/$c$; $\approx\!-52$~MeV/$c$ for
$p_z \approx 1$~GeV/$c$). Fisher 68\% $p_x$ coverage recovered from $\approx 0.05$
to $\approx 0.80$ after the fix.
Shared helper: \texttt{libs/analysis\_pdc\_reco/include/PDCFrameRotation.hh}.
\end{frame}

\begin{frame}{Study Scope}
\textbf{Physics question}
\begin{itemize}
\item Reconstruct proton target momentum from a two-point PDC track in the SAMURAI field.
\item Quantify not only the central estimate, but also the uncertainty attached to that estimate.
\end{itemize}

\vspace{0.2cm}
\textbf{This report focuses on}
\begin{itemize}
\item Runtime reconstruction backends now exposed by \texttt{reconstruct\_target\_momentum}
\item RK-based uncertainty analysis implemented in \texttt{analysis\_pdc\_reco}
\item Diagnostic comparison between RK, NN, and multidimensional regression branches
\item Confidence and credible intervals for \(p_x\), \(p_y\), \(p_z\), and \(|\vec p|\)
\end{itemize}
\end{frame}

\begin{frame}{Top-Level Reconstruction Methods}
\small
\begin{table}
\centering
\begin{tabular}{p{0.20\textwidth}p{0.28\textwidth}p{0.22\textwidth}p{0.20\textwidth}}
\toprule
Method & Input / model & Current strength & Current limitation \\
\midrule
RK fit & Field-map back-tracing from target to PDC constraints & Physics-based, interpretable, supports uncertainty propagation & Slowest branch; quality depends on initialization and fit mode \\
NN regressor & Direct map from PDC coordinates to \((p_x,p_y,p_z)\) & Best point-estimate accuracy on current truth grid & No interval output in runtime yet \\
MultiDim surrogate & Polynomial surrogate from offline regression & Fast when calibrated & Current truth-grid sample has no matched events; no interval support \\
\bottomrule
\end{tabular}
\end{table}

\vspace{0.1cm}
The current uncertainty study is therefore centered on the RK branch, because it is the only branch that now exposes Fisher, profile-likelihood, Laplace, and MCMC interval estimates.
\end{frame}

\begin{frame}{RK Reconstruction Modes}
\small
\textbf{The RK backend currently supports three fit modes}
\begin{enumerate}
\item \textbf{two-point-backprop}: infer momentum by back-propagating the PDC line to the target in the field map.
\item \textbf{fixed-target-pdc-only}: keep the target anchor fixed and fit the momentum state against the two PDC points.
\item \textbf{three-point-free}: fit both target point and momentum state as free parameters under target and PDC constraints.
\end{enumerate}

\vspace{0.2cm}
\textbf{Why this matters for uncertainty}
\begin{itemize}
\item The parameterization determines the Hessian conditioning and therefore the Fisher covariance.
\item The fit mode also changes the profile-likelihood geometry and the posterior shape.
\item In the present implementation, the enlarged legacy-run analysis uses \texttt{fixed-target-pdc-only}, because it is the most stable mode for real-style reconstructed input.
\end{itemize}
\end{frame}

\begin{frame}{What Is Written by the RK Error Analysis}
\small
\textbf{Runtime ROOT branches now available}
\begin{itemize}
\item Central fit diagnostics: status, method, iterations, raw and reduced \(\chi^2\), \(N_{\mathrm{dof}}\)
\item Fisher widths: \(\sigma_{p_x}\), \(\sigma_{p_y}\), \(\sigma_{p_z}\), \(\sigma_{|\vec p|}\)
\item Optional Laplace intervals: 68\% and 95\% lower/upper bounds for \(p_x\), \(p_y\), \(p_z\), and \(|\vec p|\)
\end{itemize}

\vspace{0.2cm}
\textbf{Offline diagnostic tool}
\begin{itemize}
\item \texttt{build/bin/analyze\_pdc\_rk\_error}
\item Outputs: \texttt{track\_summary.csv}, \texttt{profile\_samples.csv}, \texttt{bayesian\_samples.csv}, \texttt{summary.csv}, \texttt{validation\_summary.csv}, and summary plots
\item This tool re-runs the RK fit on reconstructed PDC tracks and computes frequentist and Bayesian intervals on sampled candidates
\end{itemize}
\end{frame}

\begin{frame}{Interval Definitions}
\small
\begin{table}
\centering
\begin{tabular}{p{0.17\textwidth}p{0.31\textwidth}p{0.22\textwidth}p{0.20\textwidth}}
\toprule
Estimate & Construction & Statistical meaning & Current use \\
\midrule
Fisher & Inverse local Hessian around the RK optimum & Local Gaussian approximation to the likelihood & Fast width estimate for all successful tracks \\
Profile likelihood & Scan one fit parameter, minimize others, use \(\Delta\chi^2=1,4\) crossings & Frequentist confidence interval & Asymmetry and nonlinearity diagnosis on sampled tracks \\
Laplace posterior & Gaussian posterior near the mode using posterior covariance & Bayesian credible interval & Fast Bayesian interval for all tracks if posterior is valid \\
MCMC posterior & Random-walk samples from the posterior & Bayesian credible interval with shape diagnostics & Validation on sampled tracks, not on every event \\
\bottomrule
\end{tabular}
\end{table}

\vspace{0.1cm}
\textbf{Important distinction}: Fisher and Laplace are local approximations; profile likelihood and MCMC are more robust to non-Gaussian structure but much slower.
\end{frame}

\begin{frame}{Synthetic Truth-Grid Method Comparison}
\small
\textbf{Diagnostic dataset}
\begin{itemize}
\item Existing five-point truth grid around \(p_z = 627\) MeV/c
\item One matched event per truth point in the current regression fixture
\item Use this only as a backend sanity check, not as a final performance benchmark
\end{itemize}

\begin{table}
\centering
\begin{tabular}{lrrrr}
\toprule
Backend & Matched events & Mean \(|\mathrm{bias}_{p_x}|\) & Mean \(|\mathrm{bias}_{p_y}|\) & Mean \(|\mathrm{bias}_{p_z}|\) \\
\midrule
NN & 5 & 21.04 & 15.56 & 3.10 \\
RK fixed-target & 5 & 178.35 & 11.46 & 234.39 \\
RK three-point & 5 & 651.35 & 15.77 & 190.31 \\
RK two-point & 5 & 67.99 & 1.40 & 604.92 \\
MultiDim & 0 & n/a & n/a & n/a \\
\bottomrule
\end{tabular}
\end{table}

\vspace{0.1cm}
On this fixture, the NN branch is currently the strongest point estimator, while the RK modes still need calibration or configuration tuning for this truth-grid operating point.
\end{frame}

\begin{frame}{How to Read the Truth-Grid Result}
\small
\textbf{Interpretation}
\begin{itemize}
\item The truth-grid test is very small, so the table is dominated by pointwise behaviour rather than stable statistics.
\item Even so, the pattern is useful: each RK mode fails differently.
\item \texttt{two-point-backprop} keeps transverse components relatively close, but collapses \(p_z\).
\item \texttt{fixed-target-pdc-only} is more balanced and therefore chosen as the first uncertainty-analysis mode for legacy reconstructed samples.
\item \texttt{three-point-free} is the most flexible and also the least stable on the current small truth-grid fixture.
\end{itemize}

\vspace{0.2cm}
\textbf{Practical conclusion}
\begin{itemize}
\item Use the NN branch when the immediate goal is best point prediction.
\item Use the RK branch when the goal is uncertainty-aware reconstruction and physically interpretable interval estimates.
\end{itemize}
\end{frame}

\begin{frame}{Expanded Legacy Reconstruction Sample}
\scriptsize
\textbf{Input class}
\begin{itemize}
\item Legacy reconstructed \texttt{after\_pdcana} ROOT files
\item Geometry macro: \texttt{geometry\_B115T\_pdcOptimized\_20260227.mac}
\item Magnetic field map: \texttt{180703-1,15T-3000.table}, dipole rotation \(30^\circ\)
\item RK mode: \texttt{fixed-target-pdc-only}
\end{itemize}

\vspace{0.1cm}
\textbf{Expanded run configuration}
\begin{itemize}
\item Up to 500 events per file
\item Profile scan: 11 points per parameter, 4 sampled tracks per momentum quartile
\item Bayesian validation: 2 sampled tracks per quartile, 120 kept MCMC samples after burn-in/thinning
\item Current completed output: 815 successful RK tracks, 32 profile samples, 16 Bayesian samples
\end{itemize}
\end{frame}

\begin{frame}{Reconstructed Momentum Spectrum}
\begin{center}
\IfFileExists{\LargeRunRoot/plots/reco_p_distribution.png}
  {\includegraphics[height=0.60\textheight]{\LargeRunRoot/plots/reco_p_distribution.png}}
  {\PlotOrFallback{\LargeRunRoot/plots/reco_p_distribution.png}{\SmokeRunRoot/plots/reco_p_distribution.png}{0.68}}
\end{center}

\vspace{0.1cm}
\small
The enlarged legacy sample covers a broad reconstructed proton momentum range. The median \(|\vec p|\) is \(1370.4\) MeV/c with interquartile range \(970.8\) to \(2920.8\) MeV/c.
\end{frame}

\begin{frame}{Fisher and Profile-Likelihood Widths}
\begin{columns}[T]
\column{0.50\textwidth}
\begin{center}
\PlotOrFallback{\LargeRunRoot/plots/fisher_p_width68.png}{\SmokeRunRoot/plots/fisher_p_width68.png}{0.95}
\end{center}

\column{0.50\textwidth}
\begin{center}
\PlotOrFallback{\LargeRunRoot/plots/profile_p_width68.png}{\SmokeRunRoot/plots/profile_p_width68.png}{0.95}
\end{center}
\end{columns}

\vspace{0.1cm}
\small
In the enlarged run, the median 68\% width in \(|\vec p|\) is \(48.20\) MeV/c from the Fisher estimate and \(76.91\) MeV/c from profile likelihood. The profile interval remains substantially broader, so the likelihood is still not well described by a purely local Gaussian approximation.
\end{frame}

\begin{frame}{Laplace Versus MCMC}
\begin{columns}[T]
\column{0.54\textwidth}
\begin{center}
\PlotOrFallback{\LargeRunRoot/plots/laplace_vs_mcmc_p_width68.png}{\SmokeRunRoot/plots/laplace_vs_mcmc_p_width68.png}{0.84}
\end{center}

\column{0.44\textwidth}
\scriptsize
\textbf{Enlarged run}
\begin{itemize}
\item Median Laplace 68\% width in \(|\vec p|\): \(48.14\) MeV/c
\item Median MCMC 68\% width in \(|\vec p|\): \(109.32\) MeV/c
\item Median MCMC acceptance rate: \(0.593\)
\item Median effective sample size: \(9.23\)
\end{itemize}

\vspace{0.1cm}
\textbf{Reading}
\begin{itemize}
\item Laplace tracks Fisher closely, as expected for two local Gaussian approximations.
\item MCMC is much broader on sampled tracks, which again suggests non-Gaussian posterior structure.
\end{itemize}
\end{columns}
\end{frame}

\begin{frame}{Representative Profile Curves}
\begin{center}
\PlotOrFallback{\LargeRunRoot/plots/representative_profile_curves.png}{\SmokeRunRoot/plots/representative_profile_curves.png}{0.66}
\end{center}

\vspace{0.1cm}
\scriptsize
This is the clearest frequentist diagnostic. Symmetric parabolic curves would support Fisher-like treatment. Asymmetry, flat directions, or clipped crossings explain why profile-likelihood intervals can be wider than Hessian-based widths.
\end{frame}

\begin{frame}{Confidence and Credible Interval Summary}
\small
\begin{table}
\centering
\begin{tabular}{p{0.18\textwidth}p{0.23\textwidth}p{0.22\textwidth}p{0.24\textwidth}}
\toprule
Method & Tracks covered now & 68\% interval type & Main message \\
\midrule
Fisher & All successful RK tracks & Local confidence width & Good for speed, optimistic when the likelihood is curved or asymmetric \\
Laplace & All successful RK tracks with valid posterior covariance & Local Bayesian credible interval & Similar to Fisher in the current smoke study \\
Profile likelihood & Sampled RK tracks & Frequentist confidence interval & Captures asymmetry and broader nonlocal structure \\
MCMC & Sampled RK tracks & Bayesian credible interval & Best shape information, but highest runtime cost \\
\bottomrule
\end{tabular}
\end{table}

\vspace{0.1cm}
For current PDC legacy reco data, the safest interpretation is: Fisher/Laplace for fast per-track monitoring, profile/MCMC for method validation and interval realism checks.
\end{frame}

\begin{frame}{Current Limitation on Coverage Validation}
\small
\textbf{What is missing}
\begin{itemize}
\item The legacy \texttt{after\_pdcana} files contain an \texttt{OriginalParticleData} branch without a safe dictionary in this workflow.
\item Direct truth binding from those files caused the earlier memory-corruption crash and has been disabled in the analysis tool.
\end{itemize}

\vspace{0.2cm}
\textbf{Consequence}
\begin{itemize}
\item We can compute interval sizes and compare methods on reconstructed tracks.
\item We cannot yet claim empirical 68\% or 95\% coverage on the legacy sample.
\item Coverage validation must be done on simulation files with a safe truth interface or by adding a reliable dictionary-backed legacy reader.
\end{itemize}
\end{frame}

\begin{frame}{Implementation Status}
\small
\textbf{Completed}
\begin{itemize}
\item RK runtime branches for central fit diagnostics and interval outputs
\item Offline tool for Fisher, profile-likelihood, Laplace, and MCMC comparison
\item Synthetic backend comparison CSV summaries
\item Stability fix for the legacy-file crash path caused by unsafe truth-branch binding
\end{itemize}

\vspace{0.2cm}
\textbf{Still needed}
\begin{itemize}
\item Calibrate RK modes against larger truth-labeled samples
\item Add interval support for the NN branch, for example heteroscedastic regression or quantile targets
\item Restore truth-based coverage validation for legacy-style files through a safe reader path
\end{itemize}
\end{frame}

\begin{frame}{Takeaway}
\small
\begin{itemize}
\item The PDC reconstruction framework now supports several central-estimate methods, but only the RK branch currently supports rigorous uncertainty analysis.
\item On the available legacy reconstructed sample, local Gaussian intervals are much narrower than profile-likelihood and MCMC intervals, so Gaussian-only error bars are not sufficient.
\item The NN branch remains the strongest point predictor on the current truth-grid diagnostic, but it is not yet uncertainty-aware.
\item The next technically important step is not another point-estimate tweak; it is truth-backed interval validation on a larger, stable dataset.
\end{itemize}
\end{frame}

\end{document}
